%%% Local Variables:
%%% mode: LaTeX
%%% TeX-engine: luatex
%%% TeX-master: t
%%% TeX-backend: "biber"
%%% End:
\documentclass[11pt,article,spanish]{apuntes-scr}

\universidad{Ecosistema GNU Emacs 30 / Sway / Wayland}
\curso{Manual de Referencia Oficial}

\title[Visor de PDF Minimalista con Integración SyncTeX y Wayland]{Guía Completa de Uso y Configuración de Zathura}
\author{Campos Mendoza, D. Daniel}
\date{\today\ \textperiodcentered\ \DTMcurrenttime}

\begin{document}

\maketitle
\tableofcontents
\vspace{1.5em}\hrule\vspace{1.5em}

\section{Introducción y Filosofía de Zathura}
\textbf{Zathura} es un visor de documentos altamente eficiente, modular y minimalista diseñado bajo la filosofía Unix y la ergonomía modal de Vim. Utiliza el motor de renderizado de alta fidelidad (\texttt{mupdf} o \texttt{poppler}) y está completamente integrado con el servidor gráfico Wayland y el editor GNU Emacs mediante el protocolo \textbf{SyncTeX}.

\begin{notabox}[Filosofía de Diseño]
  Zathura sigue el principio de \textit{«hacer una cosa y hacerla bien»}. No es un navegador web ni un gestor de archivos; es un visor de documentos puro que prioriza la velocidad, la legibilidad y la integración perfecta con tu flujo de trabajo en \LaTeX.
\end{notabox}

\section{Navegación y Atajos de Teclado Esenciales}
Toda la interacción en Zathura se realiza mediante el teclado, siguiendo la ergonomía modal de Vim:

\subsection{Movimiento y Desplazamiento (Scroll)}
\begin{itemize}[leftmargin=*]
    \item \textbf{\texttt{j}} / \textbf{\texttt{k}} o \textbf{$\downarrow$} / \textbf{$\uparrow$}: Desplazamiento vertical continuo por líneas.
    \item \textbf{\texttt{h}} / \textbf{\texttt{l}} o \textbf{$\leftarrow$} / \textbf{$\rightarrow$}: Desplazamiento horizontal.
    \item \textbf{\texttt{Ctrl + d}} / \textbf{\texttt{Ctrl + u}}: Avanzar / Retroceder media pantalla.
    \item \textbf{\texttt{Ctrl + f}} / \textbf{\texttt{Ctrl + b}}: Avanzar / Retroceder una página completa.
    \item \textbf{\texttt{J}} / \textbf{\texttt{K}}: Saltar directamente a la página siguiente / anterior.
    \item \textbf{\texttt{gg}} / \textbf{\texttt{G}}: Ir a la primera / última página del documento.
    \item \textbf{\texttt{[número]G}} (ej. \texttt{45G}): Saltar directamente a la página número 45.
\end{itemize}

\subsection{Visualización, Modos y Zoom}
\begin{itemize}[leftmargin=*]
    \item \textbf{\texttt{i}} (\texttt{recolor}): Alterna instantáneamente entre el \textbf{Modo Claro} (hoja blanca real) y el \textbf{Modo Oscuro} (fondo \#111111 con texto en contraste). Gracias a \texttt{recolor-keephue true}, los colores de fórmulas, gráficos y enlaces se preservan intactos.
    \item \textbf{\texttt{d}} (\texttt{toggle\_page\_mode}): Alterna entre vista de \textbf{1 página} y \textbf{2 páginas por fila} (modo lectura de libro abierto).
    \item \textbf{\texttt{a}} / \textbf{\texttt{s}}: Ajustar automáticamente la página al ancho de la ventana (\textit{fit to width}) o al tamaño completo de la pantalla.
    \item \textbf{\texttt{+}} / \textbf{\texttt{-}} / \textbf{\texttt{=}}: Aumentar zoom, disminuir zoom o restablecer zoom al 100\%.
    \item \textbf{\texttt{r}} / \textbf{\texttt{R}}: Rotar la página 90 grados / Forzar recarga manual del archivo PDF.
\end{itemize}

\subsection{Búsqueda y Marcadores (Bookmarks)}
\begin{itemize}[leftmargin=*]
    \item \textbf{\texttt{/ [texto] + Enter}}: Búsqueda difusa de palabras o expresiones en el PDF.
    \item \textbf{\texttt{n}} / \textbf{\texttt{N}}: Saltar al siguiente / anterior resultado encontrado.
    \item \textbf{\texttt{m[a-z]}} (ej. \texttt{ma}): Crear una marca rápida asociada a la letra 'a'.
    \item \textbf{\texttt{'[a-z]}} (ej. \texttt{'a}): Saltar de inmediato a la marca 'a'.
    \item \textbf{\texttt{Tab}}: Desplegar el índice del documento (Table of Contents / Bookmarks).
    \item \textbf{\texttt{y}}: Copiar el texto seleccionado directamente al portapapeles del sistema (Wayland).
\end{itemize}

\begin{controlbox}[Autoevaluación]
  ¿Qué secuencia de teclas usarías para saltar a la página 42, buscar la palabra «teorema» e ir al siguiente resultado?
\end{controlbox}

\section{Sincronización Bidireccional SyncTeX con Emacs}
La integración permite una interacción simbiótica entre el código fuente \LaTeX{} y el PDF generado:
\begin{enumerate}[leftmargin=*]
    \item \textbf{Forward Search (Emacs $\to$ Zathura)}: Al presionar \texttt{; k c} o \texttt{C-c C-a} en Emacs, tras compilar exitosamente, Zathura salta automáticamente a la página y resalta en color cian (\texttt{\#00d3d0}) la línea exacta donde está el cursor.
    \item \textbf{Inverse Search (Zathura $\to$ Emacs)}: Al hacer \textbf{\texttt{Ctrl + Clic Izquierdo}} sobre cualquier párrafo, ecuación o teorema en Zathura, el foco regresa a Emacs abriendo el archivo fuente en la línea precisa correspondiente.
\end{enumerate}

\begin{notabox}[Flujo de Trabajo SyncTeX]
  1. Edita tu \texttt{.tex} en Emacs. 2. Compila con \texttt{; k c} (Lua\LaTeX). 3. Zathura se abre mostrando la página correcta. 4. Al leer el PDF, \textbf{Ctrl+Clic} en cualquier punto te lleva de vuelta a la línea exacta en Emacs.
\end{notabox}

\section{Configuración de Apariencia y Renderizado}
El archivo reside en \texttt{\textasciitilde/.config/zathura/zathurarc} y controla cada aspecto visual y comportamental del visor:

\subsection{Integración con Emacs y SyncTeX}
\begin{lstlisting}[language=TeX, basicstyle=\ttfamily\footnotesize, frame=single, backgroundcolor=\color{gray!5}]
set synctex true
set synctex-editor-command "emacsclient -n +%{line} %{input}"
\end{lstlisting}

\subsection{Ajuste de Ventana y Desplazamiento}
\begin{itemize}[leftmargin=*]
    \item \textbf{\texttt{set adjust-open "width"}}: Abre siempre el PDF ajustado al ancho de la ventana.
    \item \textbf{\texttt{set scroll-step 40}}: Suaviza el desplazamiento línea a línea.
    \item \textbf{\texttt{set guioptions ""}}: Oculta las barras de desplazamiento para un look minimalista.
\end{itemize}

\subsection{Modo Oscuro (Recolor)}
\begin{lstlisting}[language=TeX, basicstyle=\ttfamily\footnotesize, frame=single, backgroundcolor=\color{gray!5}]
set recolor-lightcolor "#111111"
set recolor-darkcolor  "#e0e0e0"
set recolor-keephue   true
\end{lstlisting}

\begin{warningbox}[Advertencia]
  El modo \texttt{recolor} invierte los colores del PDF original. Si tu documento ya tiene un fondo oscuro nativo, desactívalo con \texttt{set recolor false} para evitar contraste excesivo.
\end{warningbox}

\subsection{Paleta de Colores de la Interfaz}
\begin{center}
\small
\begin{tabularx}{\linewidth}{l X}
\toprule
\textbf{Elemento} & \textbf{Color} \\
\midrule
\texttt{statusbar-bg} / \texttt{inputbar-bg} & \#1e1e1e (Fondo oscuro) \\
\texttt{statusbar-fg} / \texttt{inputbar-fg} & \#ffffff (Texto blanco) \\
\texttt{notification-bg} & \#d3869b (Violeta suave) \\
\texttt{notification-error-bg} & \#ff0000 (Rojo error) \\
\bottomrule
\end{tabularx}
\end{center}

\section{Atajos de Teclado Personalizados}
Además de los atajos por defecto, tu configuración añade:

\begin{center}
\small
\begin{tabularx}{\linewidth}{l X}
\toprule
\textbf{Atajo} & \textbf{Acción} \\
\midrule
\texttt{i} & Alternar Modo Oscuro / Modo Claro (\texttt{recolor}). \\
\texttt{R} & Recargar el PDF manualmente (\texttt{reload}). \\
\bottomrule
\end{tabularx}
\end{center}

\begin{notabox}[Consejo de Productividad]
  Usa \texttt{ma} \dots \texttt{mz} para marcar hasta 26 posiciones clave en un documento largo. Combínalo con la búsqueda \texttt{/} para navegar libros de 500+ páginas en segundos.
\end{notabox}

\backmatter

\end{document}
